\documentclass[11pt,border=12pt,multi=sheet]{standalone}
\usepackage[T1]{fontenc}
\usepackage{lmodern}
\usepackage{amsmath,amssymb,mathtools,mathpartir,xcolor}

\ifdefined\RulesDark
  \definecolor{ink}{HTML}{E6EDF3}
  \definecolor{muted}{HTML}{9DA7B3}
  \definecolor{accent}{HTML}{79B8FF}
\else
  \definecolor{ink}{HTML}{1F2328}
  \definecolor{muted}{HTML}{59636E}
  \definecolor{accent}{HTML}{0969DA}
\fi
\ifdefined\pdfinfoomitdate\pdfinfoomitdate=1\fi
\ifdefined\pdftrailerid\pdftrailerid{}\fi
\ifdefined\pdfsuppressptexinfo\pdfsuppressptexinfo=15\fi

\newenvironment{sheet}{%
  \begin{minipage}{160mm}\color{ink}%
  \setlength{\parindent}{0pt}\setlength{\parskip}{5pt}%
  \setlength{\abovedisplayskip}{7pt}\setlength{\belowdisplayskip}{7pt}%
}{\end{minipage}}
\newcommand{\sheettitle}[2]{%
  {\sffamily\color{accent}\small FULL $F_{<:}^{\omega}$\hfill #1 / 3}\par
  {\sffamily\bfseries\fontsize{18}{22}\selectfont #2}\par
  \vspace{2pt}{\color{muted}\hrule height .35pt}\vspace{8pt}}
\newcommand{\blocktitle}[1]{\par\vspace{5pt}{\sffamily\bfseries #1}\par}
\newcommand{\note}[1]{{\small\color{muted}#1\par}}
\newcommand{\source}[1]{\vspace{6pt}{\color{muted}\hrule height .25pt}\vspace{4pt}%
  {\scriptsize\sffamily Formal definitions: \texttt{#1}\par}}
\newcommand{\Top}{\mathrm{Top}}
\newcommand{\wf}{\mathsf{wf}}
\newcommand{\All}[3]{\forall X<:#1::#2.\,#3}
\newcommand{\TAbs}[3]{\Lambda X<:#1::#2.\,#3}
\newcommand{\ki}[3]{#1\vdash #2::#3}
\newcommand{\su}[4]{#1\vdash #2<:#3::#4}
\newcommand{\ty}[4]{#1;#2\vdash #3:#4}
\newcommand{\bound}[2]{\Gamma,X<:#1::#2}

\begin{document}

\begin{sheet}
\sheettitle{1}{Syntax and kinding}
\note{Named notation for the de Bruijn formalization. Fresh binders and
capture-avoiding substitution are understood.}

\[
\begin{array}{r@{\;}c@{\;}l}
K,L &::=& * \mid K\Rightarrow L \\[3pt]
A,B,F &::=& X \mid \Top \mid A\to B \mid \All{A}{K}{B}
                    \mid \lambda X.\,A \mid F\,A \\[3pt]
t,u &::=& x \mid \lambda x:A.\,t \mid t\,u
                 \mid \TAbs{A}{K}{t} \mid t[A]_K \\[3pt]
\Gamma &::=& \varnothing \mid \Gamma,X<:A::K \\[3pt]
\Delta &::=& \varnothing \mid \Delta,x:A
\end{array}
\]
\note{Operator-lambda kinds are carried by judgments. Term-level type
abstraction and application retain their kind annotations.}
\[
\Top_*:=\Top
\qquad
\Top_{K\Rightarrow L}:=\lambda X.\,\Top_L
\qquad
\Gamma,X::K\;:=\;\bound{\Top_K}{K}
\]

\blocktitle{Type-variable contexts}
\begin{mathpar}
\inferrule*[right=C-Empty]{\ }{\vdash\varnothing\;\wf}
\and
\inferrule*[right=C-Bind]{\vdash\Gamma\;\wf \\ \ki{\Gamma}{A}{K}}
  {\vdash\bound{A}{K}\;\wf}
\end{mathpar}

\blocktitle{Kinding}
\note{Kinding reads only the kinds in $\Gamma$, independently of bounds.}
\begin{mathpar}
\inferrule*[right=K-Var]{(X<:U::K)\in\Gamma}{\ki{\Gamma}{X}{K}}
\and
\inferrule*[right=K-Top]{\ }{\ki{\Gamma}{\Top}{*}}
\and
\inferrule*[right=K-Arrow]{\ki{\Gamma}{A}{*} \\ \ki{\Gamma}{B}{*}}
  {\ki{\Gamma}{A\to B}{*}}
\and
\inferrule*[right=K-Abs]{\ki{\Gamma,X::K}{A}{L}}
  {\ki{\Gamma}{\lambda X.\,A}{K\Rightarrow L}}
\and
\inferrule*[right=K-All]{\ki{\Gamma}{A}{K} \\ \ki{\bound{A}{K}}{B}{*}}
  {\ki{\Gamma}{\All{A}{K}{B}}{*}}
\and
\inferrule*[right=K-App]{\ki{\Gamma}{F}{K\Rightarrow L} \\ \ki{\Gamma}{A}{K}}
  {\ki{\Gamma}{F\,A}{L}}
\end{mathpar}

\blocktitle{Type conversion}
\[
(\lambda X.\,A)\,B\longrightarrow_\beta A[B/X]
\]
\note{Reduction is closed under all type constructors. Its reflexive,
symmetric, transitive closure is $\equiv_\beta$. There is no $\eta$ rule.}
\source{theories/Syntax.v; theories/HigherOrder/Subtyping.v}
\end{sheet}

\begin{sheet}
\sheettitle{2}{Declarative subtyping}
\note{The judgment $\su{\Gamma}{A}{B}{K}$ compares types at the same kind.
These rules are used with well-formed contexts.}

\begin{mathpar}
\inferrule*[right=S-Refl]{\ki{\Gamma}{A}{K}}{\su{\Gamma}{A}{A}{K}}
\and
\inferrule*[right=S-Top]{\ki{\Gamma}{A}{K}}{\su{\Gamma}{A}{\Top_K}{K}}
\and
\inferrule*[right=S-Var]{(X<:A::K)\in\Gamma \\ \ki{\Gamma}{A}{K}}
  {\su{\Gamma}{X}{A}{K}}
\and
\inferrule*[right=S-Trans]{\su{\Gamma}{A}{B}{K} \\ \su{\Gamma}{B}{C}{K}}
  {\su{\Gamma}{A}{C}{K}}
\and
\inferrule*[right=S-Arrow]{\su{\Gamma}{A'}{A}{*} \\ \su{\Gamma}{B}{B'}{*}}
  {\su{\Gamma}{A\to B}{A'\to B'}{*}}
\and
\inferrule*[right=S-All]{\su{\Gamma}{A'}{A}{K} \\ \su{\bound{A'}{K}}{B}{B'}{*}}
  {\su{\Gamma}{(\All{A}{K}{B})}{(\All{A'}{K}{B'})}{*}}
\and
\inferrule*[right=S-Abs]{\su{\bound{\Top_K}{K}}{A}{B}{L}}
  {\su{\Gamma}{\lambda X.\,A}{\lambda X.\,B}{K\Rightarrow L}}
\and
\inferrule*[right=S-App]{\su{\Gamma}{F}{F'}{K\Rightarrow L} \\ \ki{\Gamma}{C}{K}}
  {\su{\Gamma}{F\,C}{F'\,C}{L}}
\and
\inferrule*[right=S-Conv]{\ki{\Gamma}{A}{K} \\ \ki{\Gamma}{B}{K} \\ A\equiv_\beta B}
  {\su{\Gamma}{A}{B}{K}}
\end{mathpar}

\blocktitle{Full quantification and pointwise operators}
\note{S-All permits different bounds. The bound is contravariant, and
the bodies are compared under the right-hand bound $A'$. S-Abs uses an
unbounded operator parameter. S-App uses the same argument $C$ on both
sides; no general covariance rule for operator arguments is present.}
\note{S-Trans permits arbitrary well-kinded intermediate types. The erasure
used in the conservativity proof is not an object-language conversion rule.}
\source{theories/HigherOrder/Subtyping.v}
\end{sheet}

\begin{sheet}
\sheettitle{3}{Explicitly annotated typing}
\note{$\Gamma$ contains type-variable bounds; $\Delta$ contains term-variable
types. The result type in $\ty{\Gamma}{\Delta}{t}{A}$ has kind $*$.}

\begin{mathpar}
\inferrule*[right=T-Var]{(x:A)\in\Delta \\ \ki{\Gamma}{A}{*}}
  {\ty{\Gamma}{\Delta}{x}{A}}
\and
\inferrule*[right=T-Abs]{\ki{\Gamma}{A}{*} \\ \ty{\Gamma}{\Delta,x:A}{t}{B}}
  {\ty{\Gamma}{\Delta}{\lambda x:A.\,t}{A\to B}}
\and
\inferrule*[right=T-App]{\ty{\Gamma}{\Delta}{t}{A\to B} \\ \ty{\Gamma}{\Delta}{u}{A}}
  {\ty{\Gamma}{\Delta}{t\,u}{B}}
\and
\inferrule*[right=T-TAbs]{\ki{\Gamma}{A}{K} \\ \ty{\bound{A}{K}}{\Delta}{t}{B}}
  {\ty{\Gamma}{\Delta}{\TAbs{A}{K}{t}}{\All{A}{K}{B}}}
\and
\inferrule*[right=T-TApp]{\ty{\Gamma}{\Delta}{t}{\All{A}{K}{B}} \\ \su{\Gamma}{C}{A}{K}}
  {\ty{\Gamma}{\Delta}{t[C]_K}{B[C/X]}}
\and
\inferrule*[right=T-Sub]{\ty{\Gamma}{\Delta}{t}{A} \\ \su{\Gamma}{A}{B}{*}}
  {\ty{\Gamma}{\Delta}{t}{B}}
\end{mathpar}

\note{In T-TAbs, $X$ is fresh for the existing contexts. The named notation
keeps $\Delta$ unchanged; the de Bruijn implementation shifts its stored types.}

\blocktitle{The old fragment}
\[
S,T ::= X\mid\Top\mid S\to T\mid\All{S}{*}{T}
\]
\note{Old contexts bind only base-kind type variables with old bounds.
Old terms use base-kind type abstraction and application, with old lambda
annotations and old type arguments. The proved conservativity results
compare judgments whose contexts, terms, and endpoints satisfy these
restrictions.}
\note{The formalization specifies kinding, subtyping, and typing. It does
not define an operational semantics for terms.}
\source{theories/Terms.v; theories/HigherOrder/Typing.v}
\end{sheet}

\end{document}
